\documentclass[../../../include/open-logic-section]{subfiles}

\begin{document}

\olfileid{sth}{replacement}{intro}
\olsection{పరిచయం}

$\ZF$,~$\Z$ల మధ్య తేడాకు కారణమైన స్వయంసిద్ధ పథకం ప్రతిస్థాపన. \crefrange{sth:ordinals::chap}{sth:spine::chap} అంతటా దాన్ని వాడుకున్నాం. ఈ అధ్యాయంలో చివరకు ఈ ప్రశ్నను పరిశీలిస్తాం: ప్రతిస్థాపన సమర్థితమేనా?

ప్రశ్నను పదును పెట్టడానికి, ప్రతిస్థాపన నిజానికి చాలా \emph{బలమైనది} అని గమనించాలి. ఈ అధ్యాయంలో (మళ్లీ \olref[sth][card-arithmetic][fix]{sec}లో) అది ఎంత బలమైనదో తెలుస్తుంది. అందువల్ల దాని సమర్థన నిజంగానే అవసరమని తెలుస్తుంది.

రెండు రకాల సమర్థనలను చర్చిస్తాం. స్థూలంగా: \emph{బాహ్య} సమర్థన ఒక స్వయంసిద్ధ సూత్రాన్ని దాని ఫలితాల ఆధారంగా సమర్థించే ప్రయత్నం; \emph{అంతర్గత} సమర్థన సంబంధిత గణిత భావనలే దానిని సమర్థిస్తాయని చూపించే ప్రయత్నం. దీని అర్థం ఈ అధ్యాయంలో మరింత స్పష్టమవుతుంది; అయితే ఇది పెద్ద చర్చలో చిన్న భాగమే. మరింత కోసం ముఖ్యంగా \citeauthor{Maddy1988a} (\citeyear{Maddy1988a}, \citeyear{Maddy1988b}) చూడండి.

\end{document}
